\section{Conclusion}
\label{sec:conclusion}

In this work, we have deconstructed the growing architectural complexity of dynamic model merging through the lens of Occam's razor. By investigating the reported failure of classical linear routing—specifically on high-variance, challenging datasets like SVHN—we demonstrated that representation collapse is not an inherent limitation of linear networks. Instead, it is a standard, addressable overfitting and variance phenomenon caused by unregularized optimization on small calibration sets.

To address this, we proposed \textbf{Robust Linear Routing (RLR)}, a minimalist dynamic model merging framework. RLR retains the simple, mathematically transparent structure of a single 768-parameter linear routing layer, but stabilizes its routing outputs via standard, classical regularizations: $L_2$ weight decay and Softmax Temperature scaling. Our multi-seed experiments on a 4-task Vision Transformer benchmark demonstrate that both classical unregularized linear routing and RLR achieve outstanding, stable multi-task convergence across random calibration sets, outperforming complex baselines by large margins. Concurrently, RLR acts as a stabilizer, securing a robust $91.20\% \pm 1.84\%$ SVHN accuracy and demonstrating superior resilience under heterogeneous mixed-task evaluation streams compared to the unregularized classical baseline.

Our work serves as a reminder that before introducing convoluted, multi-stage paradigms—such as quantum-inspired wavefunctions—we must first thoroughly regularize and understand classical neural network baselines. RLR achieves outstanding parameter fusion with zero runtime overhead, absolute simplicity, and high mathematical readability. While our evaluation focused on standard compact Vision Transformers to directly deconstruct prior work, scaling classical regularized gating to high-dimensional spaces---such as Large Language Models (LLMs) with billions of parameters---remains a highly promising avenue. Specifically, we outline concrete pathways for scaling RLR to modern LLM architectures: 
(1) \textbf{Sequence-Level Pooled Routing Signals over LoRA Experts:} Instead of processing high-frequency token representations which would cause massive runtime latency, the routing network can process sequence-level pooled hidden states to dynamically blend low-rank adapters (LoRA; \cite{Hu2022}). Formally, let $\mathbf{h} \in \mathbb{R}^d$ represent the globally average-pooled hidden state of the first transformer block (or the hidden state corresponding to the initial \texttt{[BOS]} token). RLR maps this representation to blending coefficients $\boldsymbol{\lambda} \in \mathbb{R}^M$ for $M$ expert LoRA adapters via:
\begin{equation}
\boldsymbol{\lambda} = \mathrm{softmax}\left( \frac{\mathbf{W}_r \mathbf{h} + \mathbf{b}_r}{T} \right)
\end{equation}
where $\mathbf{W}_r \in \mathbb{R}^{M \times d}$ and $\mathbf{b}_r \in \mathbb{R}^M$ are the regularized routing weights and biases, and $T$ is the temperature parameter. The merged adapter update is then dynamically computed as $\Delta \mathbf{W}_{\text{merged}} = \sum_{m=1}^M \lambda_m \Delta \mathbf{W}_m$, which introduces negligible runtime and memory overhead.
(2) \textbf{High-Dimensional Weight Spaces and Linear Mode Connectivity:} Preliminary insights on high-dimensional weight spaces suggest that large foundation models fine-tuned from a shared pre-trained initialization exhibit strong linear mode connectivity~\cite{Wortsman2022}. This structural alignment implies that the low-dimensional, smooth routing trajectory guaranteed by RLR's $L_2$ decay and temperature scaling is highly compatible with the linear landscapes of billion-parameter models, suggesting that its stabilizing effects will scale seamlessly to massive architectures.

We hope this study inspires model merging researchers to favor elegant, minimalist designs over needlessly complex architectures, reinforcing the enduring value of simplicity in deep learning.
